\documentclass[10pt,a4paper]{amsart}
\usepackage[margin=19mm]{geometry}
\usepackage{amsmath,amssymb,mathtools}
\usepackage[scr=rsfs,cal=euler]{mathalfa}
\usepackage{xfrac}
\usepackage[unicode,hidelinks,bookmarks=false]{hyperref}
\newcommand{\ie}{i.e.\ }
\newcommand{\cf}{cf.\ }
\newcommand{\resp}{respectively\ }
\newcommand{\Subsection}{\subsection}
\providecommand{\nobreakdash}{\nobreak\hbox{-}}
\setlength{\emergencystretch}{2em}
\multlinegap0pt
\usepackage[shortlabels]{enumitem}
\usepackage{amssymb}
\usepackage{braket}
\usepackage{breqn}
\usepackage{dsfont}
\usepackage{mathtools}
\usepackage{float}
\usepackage{xcolor}
\usepackage{tikz}
\newtheorem{assumption}{Assumption}
\newtheorem{proposition}{Proposition}
\newtheorem{lemma}{Lemma}
\newtheorem{theorem}{Theorem}
\newcommand{\JMS}[1]{{\color{orange} {\bf J. Siktar:} [#1] }}
\def\Om{\Omega}
\def\R{\mathbb{R}}
\def\al{\alpha}
\def\be{\beta}
\def\de{\delta}
\def\ep{\epsilon}
\def\grad{\nabla}
\def\sig{\sigma}
\def\th{\theta}
\title{A Mountain-Pass Algorithm for Nonlocal Problems with Super-quadratic Nonlinearities}
\author{Loic Cappanera, Gabriela Jaramillo, Joshua M. Siktar}
\date{Source: arXiv:2605.23106v1 (CC BY 4.0)}
\begin{document}
\maketitle

\noindent\emph{Source and licence.} A Mountain-Pass Algorithm for Nonlocal Problems with Super-quadratic Nonlinearities. Version arXiv:2605.23106v1 (CC BY 4.0), distributed by arXiv under the Creative Commons Attribution 4.0 International licence, http://creativecommons.org/licenses/by/4.0/, as recorded in the arXiv licence metadata for this exact version and shown on the abstract page https://arxiv.org/abs/2605.23106v1. Full licence notice: \texttt{licenses/arxiv-2605.23106-CC-BY-4.0.txt}.\par\smallskip

\noindent\textit{Editorial scope.} Counted unit: the article's lemma l:min\_max (source0.tex lines 1263--1266). The article's complete original proof of it is delivered with it (source0.tex lines 1268--1323, one \texttt{proof} environment of 56 lines). No mathematics is written, repaired or re-derived: the statement and the proof are the article's own lines, in the article's own notation, and no retained line is substituted or re-flowed.  Retained uncounted as the source-local context the proof needs: lines 320--325; lines 813--844; lines 956--965; lines 968--974; lines 1041--1043; lines 1047--1051; lines 1233--1235; lines 1239--1247; lines 1241--1243.  Nothing was omitted from any retained span, so the package carries no omission record.  No retained line contains a citation, so the wrapper carries no bibliography and no bibliography entry.  No retained line contains a figure environment, an image-inclusion command or drawing code, so no image is delivered and the package declares no asset dependency.  Editorial formatting only, in addition: an AMS \texttt{amsart} preamble that restates the article's own declarations and macros used by the retained lines, namely the environments \texttt{assumption}, \texttt{proposition}, \texttt{lemma}, \texttt{theorem} on separate counters, together with the standard packages those lines need.  The retained environments print in this document as Assumption 1, Proposition 1, Lemma 1, Theorem 1 in order of appearance, whereas the article prints the same items with the numbering of its own full document.  The complete native source of the article as distributed in the arXiv e-print for this exact version is bundled unchanged as \texttt{sources/201184/source0.tex}.
\begin{equation}\label{e:nonlocal_eq}
\begin{array}{c c c c }
-\mathcal{L} u(x) \ & = \ f(x,u) & \quad  x \in & \Omega \subset \mathbb{R}^n \\ %\quad n=2,\\
%u & = 0 & \quad x \in & \Omega^c = \mathbb{R}^n \setminus \Omega
\end{array}
\end{equation}

\begin{assumption}[Assumptions on data]\label{Assump: AdmissibleControls}
     Let $a_1, a_2 > 0$, $\mu > 2$, $r \in (n + 1, \infty)$ be fixed constants, and let $\al \in (1, 2^*/2)$ be a fixed exponent where $2^* := \frac{2n}{n - 2}$ when $n \geq 3$ and $2^* := \infty$ when $n \leq 2$. 
     %Also let $C: (0, \infty) \rightarrow (0, \infty)$ be fixed. 
     We consider right-hand side datum $f: \Om \times \R \rightarrow \R$ satisfying the following conditions:
    \begin{enumerate}[label=\textbf{A\arabic*}]
    % \item \label{A1} We have the inequality
    %  $\|f\|_{W^{1, r}(\Om \times (-M, M))} \ \leq \ C(M)$ for all $M > 0$. 
        \item \label{A2} We have that 
        \begin{equation}\label{Eq: fGrowth}
            |f(x, t)| \ \leq \ a_1 + a_2|t|^{\al}
        \end{equation}
        for all $x \in \Om$ and $t \in \R$.
        \item\label{A3} We have the limit
        \begin{equation}\label{Eq: uniformLimit}
            \lim_{t \rightarrow 0}\frac{f(x, t)}{t} \ = \ 0
        \end{equation}
        uniformly over $x \in \Om$.
        \item\label{A4} There exists $\th \geq 1$ so that the following holds independently of $f$: if we define 
        \begin{equation}\label{Eq: Gantideriv}
            F(x, t) \ := \ \int^{t}_{0}f(x, \tau)d\tau
        \end{equation}
        then $\th\mathcal{F}(x, t) \geq \mathcal{F}(x, \sig t)$ for $(x, t) \in \R^n \times \R$ and $\sig \in [0, 1]$, where we have defined
        \begin{equation}\label{Eq: ScrFDef}
            \mathcal{F}(x, t) \ := \ tf(x, t) - \mu F(x, t).
        \end{equation}
        \item\label{A5} We have the limit
        \begin{equation}\label{Eq: uniformLimitInfinity}
            \lim_{|t| \rightarrow \infty}\frac{f(x, t)}{t} \ = \ \pm \infty
        \end{equation}
        uniformly over $x \in \Om$, and at least one of the limits $t \rightarrow \pm \infty$ is such that \eqref{Eq: uniformLimitInfinity} equals $+\infty$.
    \end{enumerate}
\end{assumption}

\begin{proposition}\label{Prop: fB}
    Let $f: \Om \times \R \rightarrow \R$ be a Carathéodory function satisfying \ref{A2} and \ref{A3}, then for any $\ep > 0$ there exists $\de(\ep) > 0$ so that for a.e. $x \in \Om$ and all $t \in \R$,
    \begin{equation}\label{Eq: fB}
        |f(x, t)| \ \leq \ 2\ep|t| + (\al + 1)\de(\ep)|t|^{\al}
    \end{equation}
        and
    \begin{equation}\label{Eq: FB}
        |F(x, t)| \ \leq \ \ep|t|^2 + \de(\ep)|t|^{\al + 1}.
    \end{equation}
    \end{proposition}

\begin{lemma}\label{Lem: theRingLemma}
    There exist constants $\rho, \be > 0$ such that, for all $u \in X$ with 
    $\| u\|_{H^1(\Omega)} = \rho$,
    %$\|\grad u\|_{L^2(\R^n)} = \rho$, 
    we have $I[u] > \be$.
  %  \JMS{Notice that $\|\grad \cdot\|_{L^2(\R^n)}$ is a norm on $X^1_D$ and is arguably the most convenient one to use here.}
\end{lemma}

  \begin{equation}\label{e:F_even}
    \lim_{|t| \rightarrow \infty}\frac{F(x, t)}{t^2} \ = \ +\infty
    \end{equation}

 \begin{lemma}\label{Lem: coercivityLemH^1_v2}
Let $\rho > 0$ be as defined in Lemma \ref{Lem: theRingLemma} and 
assume $F(x,t)$ satisfies \eqref{e:F_even}. Then, there is a constant $\rho_2>\rho$
such that for all $e \in X$ with norm $\| e\|_{H^1(\Omega)} > \rho_2$ we have that $I[e] \leq 0$.
\end{lemma}

\begin{equation}\label{e:energy_repeat}
 I[u] \ = \ \frac{1}{2} B[u,u] - \int_\Omega F(x,u) \;dx.
 \end{equation}

\begin{theorem}\label{t:existence_2}
Let $I: X \longrightarrow \R$, given in \eqref{e:energy_repeat}, be the energy associated with the nonlocal problem \eqref{e:nonlocal_eq}. If we further assume that the nonlinear function $F(x,t)$ satisfies
\begin{equation}\label{a:F}
 \lim_{|t| \to \infty} \frac{F(x,t)}{t^2} \ = \ \infty
 \end{equation}
uniformly for a.e. $ x\in \Omega$, then
\[ \bar{c} \ = \ \inf_{w \in X, w\neq 0} \max_{t\geq 0} I[tw]\]
is a critical value of $I[\cdot]$.
\end{theorem}

\begin{equation}
 \lim_{|t| \to \infty} \frac{F(x,t)}{t^2} \ = \ \infty
 \end{equation}

\begin{lemma}\label{l:min_max}
Let $I: X \longrightarrow \R$, satisfy the assumptions of Theorem \ref{t:existence_2}. Then,
\[ \bar{c} \ = \ \inf_{w \in X, w\neq 0} \max_{t\geq 0} I[tw] \ = \ \min_{w \in X, w\neq 0} \max_{t\geq 0} I[tw].  \]
\end{lemma}

\begin{proof}
Let $w \in X$ with $\|w \|_{H^1(\Omega)} =1$, and define
\[ g_w(t) \ = \ I[tw] \ = \ t^2 B[w,w] - \int_\Omega F(x,tw) \;dx.\]
Notice that $g: \R_+ \to \R$ is a smooth function of $t$. 
We want to show that $g$ attains its maximum on $\R_+$.

First, Proposition \ref{Prop: fB} implies that $g_w(0) =0$. Next, Lemma \ref{Lem: theRingLemma} shows there are numbers $\rho,\beta >0$ such that for $t= \rho$ the function $g_w(t) >\beta$, independently of $w$.
Since $F(x,t)$ satisfies property \ref{a:F}, from Lemma \ref{Lem: coercivityLemH^1_v2} we have that there exists a constant $\rho_2> \rho$ such that for all $t\geq \rho_2$ the function $g_w(t) <0$, independently of $w$. It follows that $g_w(t)$ attains its maximum for some value $t^* \in (0, \rho_2)$.


Let $G(u) = g_w(t^*)$, with $u := t^* w$ and notice that
\[ \bar{c} \ = \ \inf_{\substack{w \in X,\\ w\neq 0}} \max_{t\geq 0} I(tw) \ = \ \inf_{u \in K } G(u)\]
where $K$ is the closed ball of radius $\rho_2$ in $X$, i.e.,
\[ K \ = \ \{ u \in X: \| u\|_{H^1(\Omega)} \leq \rho_2\}.\]


{
From the compact Sobolev embedding $H^1(\Omega) \subset \subset L^p(\Omega)$, valid for any $p \in (2, 2^*)$ with $2^* =  2n/(n-2)$, we see in addition that $K$ is a compact subset of $L^p(\Omega)$. Since $F$ satisfies Assumption \ref{Assump: AdmissibleControls} and picking $ p>\alpha +1>2$, it follows that $G: L^p(\Omega) \longrightarrow \R$ is a bounded functional. Indeed, we can write
\[ |G(u)| \ \leq \ |B[u,u] | +  \int_\Omega \left| F(x,u) \right| \;dx\]
and bound first
\begin{align*}
B[u,u] \ & = \ (- \mathcal{L} u, u)_{L^2(\Omega)} \\
| B[u,u] | & \ \leq \ \|\mathcal{L}u \|_{L^2(\Om)} \| u\|_{L^2(\Omega)} \\
| B[u,u] | & \ \leq \ \|\mathcal{L}\| \|u \|^2_{L^2(\Om)} \\
| B[u,u] | & \ \leq \ C \|\mathcal{L}\| \|u \|^2_{L^p(\Om)},
\end{align*}
where the last inequality follows from the embedding $L^p(\Omega) \subset L^r(\Omega)$, valid for bounded domains $\Omega$ when $p>r$, and the fact that the operator $\mathcal{L}:L^2(\Omega) \longrightarrow L^2(\Omega)$ is bounded.

Next, from Proposition \ref{Prop: fB}, we have that for any $\ep >0$ there is a $\delta(\ep) >0$ such that 
\[ \left| F(x,u) \right| \ \leq \ \ep |u|^2 + \delta(\ep) |u|^{\alpha+1}.\]
As a result,
\begin{align*}
\int_\Omega \left| F(x,u) \right| \;dx & \ \leq \ \int_\Omega \ep |u|^2 + \delta(\ep) |u|^{\alpha+1} \;dx\\
\int_\Omega \left| F(x,u) \right| \;dx & \ \leq \ \ep \|u\|_{L^2(\Omega)}^2 + \delta(\ep)  \|u\|_{L^{\alpha+1}(\Omega)}^{\alpha+1} \\
\int_\Omega \left| F(x,u) \right| \;dx & \ \leq \ C( \ep \|u\|_{L^p(\Omega)}^2 + \delta(\ep)  \|u\|_{L^{p}(\Omega)}^{\alpha+1} ) ,
\end{align*}
where again the last inequality follows from selecting $p> \alpha +1$ and using the embedding $L^p(\Omega) \subset L^{\alpha+1}(\Omega)$.


We therefore conclude that the functional $G:L^p(\Omega) \longrightarrow \R$  attains its infimum in $K$, and there is a $\bar{u} \in K$ such that 
\[  \inf_{u \in K} G(u) \ = \ \bar{c} \ = \ G(\bar{u}) \ = \  \min_{u \in K } G(u).\] 
}









%From the compact Sobolev embedding $H^1(\Omega) \subset \subset L^p(\Omega)$ with $2< p< 2^* = 2n/(n-2)$, we see in addition that $K$ is a compact subset of $L^p(\Omega)$. 

%Since $F$ satisfies Assumption \ref{Assump: AdmissibleControls} and $ 2<p<2^*$, it follows that $G: L^p(\Omega) \longrightarrow \R$ is a bounded functional. Consequently it attains its infimum in $K$ and there is a $\bar{u} \in K$ such that 
%\[  \inf_{u \in K} G(u) = \bar{c}  = G(\bar{u})  =  \min_{u \in K } G(u).\] 
\end{proof}

\end{document}
